\chapter{系统设计}
\label{chap:design}

\section{产品形态与非目标}

PMFS v3 的产品形态是：\textbf{保留 virtio 根，经固定合成挂载点 \texttt{/pm} 提供受限文件子系统}。用户可用普通 \texttt{open}/\texttt{read}/\texttt{write}/\texttt{lseek}/\texttt{fsync}/\texttt{getdents} 等访问 PM 文件；并保留专用 KV/stream 接口。\textbf{不}替换根 FS，\textbf{不}做动态 mount table 或完整 POSIX。

非目标：上游通用 NVDIMM/NFIT/热插拔；无平台假设的物理掉电证明；符号链接、权限、\texttt{mmap}、跨后端 rename、打开后删除最后链接；全设备 GC/硬件磨损均衡。

\section{地址空间与后端}

\begin{lstlisting}[caption={地址布局示意（MAP / 实验 NVDIMM）},label={lst:layout}]
KERNBASE ---------- 0x80000000
  kernel + user（易失 DRAM，典型 128MiB）
PHYSTOP  ---------- KERNBASE + 128MiB
PMEMBASE ---------- PHYSTOP（如 0x88000000）
  PMFS 窗口：双槽 checkpoint / 元数据 / run
PMEMEND  ---------- PMEMBASE + 32MiB（名义；MAP 受 DTB 截断）
\end{lstlisting}

\begin{table}[htbp]
  \centering
  \caption{物理后端与可验证范围}
  \label{tab:backends}
  \begin{tabular}{p{2.2cm}p{5.0cm}p{5.2cm}}
    \toprule
    模式 & 含义 & 不可据此声称 \\
    \midrule
    EMUL & 8\,MiB 静态易失缓冲 & 跨启动保留、硬件持久 \\
    普通 MAP & \texttt{PHYSTOP} 后 guest RAM & 独立 NVDIMM 设备 \\
    共享 RAM 文件 & file-backend，便于杀进程重启 & 宿主掉电、独立设备 \\
    实验 NVDIMM & QEMU 9.2.0 补丁冷插拔 \texttt{nvdimm}+私有 FDT
      & 上游通用 virt、NFIT、物理掉电 \\
    Zicbom（可选） & 发出 \texttt{CBO.CLEAN} & TCG 下的真实缓存写回 \\
    \bottomrule
  \end{tabular}
\end{table}

MAP 在窗口不足或设备缺失时\textbf{拒绝挂载}，不静默降级为 EMUL。

\section{软件分层与对象图}

自上而下：VFS 逐组件遍历（原生 inode 与 PM inode pin）$\rightarrow$ \texttt{pmfile}（共享工作缓冲、COW 元数据）$\rightarrow$ \texttt{pmstream}/\texttt{pmcdc} $\rightarrow$ \texttt{pmalgo}/\texttt{pmindex} $\rightarrow$ \texttt{pmfs}（双槽 checkpoint、可达性、页复用）$\rightarrow$ \texttt{pmdev}/\texttt{pmem}。

文件内容以 KV chunk（extent）承载；元数据为不可变快照页。容量边界（量级）：约 40 inode/dentry、每文件 4 extent、文件 $\le$1024\,B。

\section{崩溃可见点与 POSIX 差集}
\label{sec:semantics}

\begin{table}[htbp]
  \centering
  \caption{提交可见性（软件协议）}
  \label{tab:visibility}
  \begin{tabular}{p{3.4cm}p{2.4cm}p{6.5cm}}
    \toprule
    操作 & 故障后承诺 & 说明 \\
    \midrule
    未成功 \texttt{fsync} 的写 & \textbf{不承诺}保留 & 工作缓冲/memtable \\
    成功 \texttt{fsync} & 完整新版本可读回 & 数据+元数据+提交词 \\
    空白/v2/校验损坏 & 拒绝挂载 & 不自动 format \\
    单槽损坏 & 可回退另一槽 & 双槽均坏则失败 \\
    \bottomrule
  \end{tabular}
\end{table}

差集包括：动态挂载、完整跨目录/跨后端 rename、符号链接、权限、\texttt{mmap}、多写者线性化证明等。硬链接、\texttt{lseek}/\texttt{O\_APPEND}、PM 内子树 rename、\texttt{getdents} 已支持（见实现章）。

\section{刷写与 checkpoint 发布顺序}

持久化请求语义依赖平台：默认 \texttt{pmem\_flush} 为按行 touch+\texttt{fence}（可观测、可排序）；可选 Zicbom 发出 \texttt{CBO.CLEAN}。一次根发布的设计顺序为：

\begin{enumerate}[leftmargin=2em]
  \item 追加 run / 分配 COW 元数据页；对各页请求 persist + fence；
  \item 非当前槽：清提交词并 persist；写入新头（未提交）并 persist；
  \item \textbf{一次 8\,B store}写入提交状态与校验 $\rightarrow$ persist + fence；再更新内存 epoch/槽。
\end{enumerate}

该协议还依赖：屏障顺序、提交词故障原子性、映射有效、校验未碰撞等假设。QEMU TCG 不模拟缓存持久域，故正文将「掉电安全」表述为\textbf{声明假设下的协议正确性}，并以进程重启 oracle 验证。

\section{维护、CDC 与评测意图}

双槽可达性保护两代根；双槽均不可达页可复用。显式 \texttt{pmcompact} 做全链 COW 压缩，近满或缺暂存空间可失败。普通写优先 CDC，超预算则固定分块；\textbf{无内容去重}。

评测分两平面：\textbf{架构面}——相对 xv6 的同工作量延迟/吞吐（持久契约不等价）；\textbf{公平面}——同 PMFS 内软件请求/策略消融。禁止混读。
